\subsection{International comparisons by repository domain}
\label{sec:sector_sc_results}

The international extension complements the before--after sector comparison with a separate comparison-country counterfactual for each fixed pre-treatment repository-domain cohort. The main historical-text rule assigns 2,163 of the 10,241 Bangladesh accounts (21.1 percent) to a domain. Of these, 1,822 accounts belong to 8 named domains that pass the support rule; the remaining eligible accounts are either in unsupported named groups or the separately reported unclassified group. Coverage by country and the full support screen are reported in \cref{tab:sector_coverage,tab:sector_support}. The complete country-by-country appendix (\cref{app:country_atlas}) displays all 20 groups, including unclassified accounts, across all 15 countries: 300 panels, of which 45 have no assigned users and 144 contain fewer than 30. These cells remain visible even when they cannot support estimation. A country-domain denominator of zero is not a zero activity rate, and small cells can produce large daily percentage movements.

\begin{table}[htbp]
\centering
\small
\caption{International synthetic controls by pre-treatment repository domain}
\label{tab:sector_sc}
\begin{threeparttable}
\begin{tabular}{@{}p{0.23\textwidth}rrrrrr@{}}
\toprule
Domain & $N_{BD,s}$ & $J_s$ & ATT (pp) & Rel. (\%) & Pre RMSPE & Rank \\
\midrule
Retail & 368 & 11 & -15.69 & -91.8 & 2.91 & 0.083 \\
Transport & 44 & 5 & -11.31 & -96.8 & 5.85 & 0.167 \\
Information & 518 & 14 & -12.74 & -90.4 & 2.63 & 0.067 \\
Finance & 90 & 9 & -12.28 & -85.8 & 4.04 & 0.100 \\
Professional services & 151 & 11 & -16.13 & -89.6 & 3.91 & 0.083 \\
Education & 424 & 13 & -15.04 & -88.2 & 4.91 & 0.071 \\
Health & 155 & 5 & -16.84 & -91.8 & 7.21 & 0.167 \\
Arts and recreation & 72 & 7 & -12.33 & -94.7 & 5.31 & 0.250 \\
Unclassified & 8,078 & 14 & -12.77 & -85.8 & 3.41 & 0.067 \\
\bottomrule
\end{tabular}
\begin{tablenotes}[flushleft]
\footnotesize
\item Main assignment: historical README and score at least 0.80. All supported domains and the unclassified group are shown. Outcome: primary public activity anywhere on GitHub among users assigned before treatment. ATT averages July 18--23. Rank compares absolute ATT with country placebos; its minimum depends on donor count. These are conditional comparative estimates, not effects on total industrial output or evidence that domain effects differ statistically.
\end{tablenotes}
\end{threeparttable}
\end{table}


Across the 8 supported named domains, estimated July 18--23 changes range from -16.84 to -11.31 percentage points. \Cref{fig:advisor_sector_sc} shows every supported estimate and the range obtained when one donor is omitted. These ranges are sensitivity summaries, not confidence intervals. \Cref{tab:sector_fit} reports pre-fit and scale-adjusted placebo ranks, and \cref{tab:sector_sc_sensitivity} retains all classification-threshold comparisons. Smaller domain samples and fewer eligible donors make the reference distributions coarser than in the national analysis. No domain ranking is treated as a statistically established ranking of causal exposure.

Retail- and Information-associated cohorts provide relatively closer pre-period fits: their RMSPEs are 2.91 and 2.63 percentage points, compared with estimated shutdown gaps of -15.69 and -12.74. Health has a larger point estimate (-16.84) but also a pre-period RMSPE of 7.21; Education's earlier-fit holdout error is 9.12 points. These diagnostics weaken a comparison of the largest point estimates. Every leave-one-donor-out estimate remains negative, but this robustness does not resolve imperfect fit, classifier error or common exposure to unrest.

\begin{figure}[htbp]
\centering
\caption{Synthetic-control gaps across pre-treatment repository domains}
\label{fig:advisor_sector_sc}
\includegraphics[width=\textwidth]{fig_advisor_sector_sc.pdf}
\fignote{Points show the July 18--23 ATT for all supported groups. Horizontal segments span the estimates from omitting each donor country in turn; they are not sampling confidence intervals. Assignment uses pre-shutdown activity and historical README classification at score 0.80. The outcome follows assigned users' public activity across repositories.}
\end{figure}

The supported named groups comprise 17.8 percent of eligible Bangladesh accounts. Weighting their estimates by their shares of the national eligible sample gives a contribution of -2.57 percentage points. This is a partial decomposition for the covered groups, not a national treatment effect: unsupported groups have no assigned zero effect, and the unclassified group is substantial. Moreover, domain-specific weights differ from the national weights. The national headline therefore remains the original $-$13.26 percentage points.

These comparisons strengthen the design relative to the older within-Bangladesh recovery indices because they introduce untreated-country counterfactuals and freeze group membership before the shock. They still depend on country-domain comparability, imperfect text classification and the absence of coincident group-specific shocks. They concern public activity among users associated with a repository domain; they do not establish effects on firms, employment sectors or industrial output. The separate recovery indices and their broader Bangladesh sample are preserved in the appendix (\cref{sec:sector_results}).
